\chapter{Conclusion et perspectives}
Ce projet a permis de concevoir et developper une plateforme de gestion des projets de stage pour la Digital Factory de Royal Air Maroc. La Digital Factory designe le departement d'accueil et de developpement, tandis que la plateforme correspond a la solution realisee durant le stage. Elle centralise les principaux processus lies au suivi des projets de stage : gestion des utilisateurs, affectations, taches, soumissions, commentaires et suivi de l'avancement. L'integration d'une generation assistee par IA offre un gain de temps a la creation des projets, tout en preservant la decision humaine grace a une etape de validation par le SuperAdmin.

Le choix d'une architecture Spring Boot, React avec TypeScript et PostgreSQL apporte une base structurée, maintenable et adaptée à une application métier. La chaîne Jenkins enrichie de contrôles DevSecOps contribue à renforcer la qualité et la sécurité du processus de livraison.

Plusieurs évolutions peuvent être envisagées : ajout de notifications par courriel, tableaux de bord analytiques, gestion fine des échéances, recherche avancée, export de rapports de suivi, amélioration de la collaboration en temps réel et enrichissement de l'assistance IA. Ces perspectives pourront être priorisées selon les besoins opérationnels de la Digital Factory.
